% Parameter and strategy record for the PMMA-RSF production campaign.
% The source of truth for numerical values is cases/rsf_0117...rsf_0126.
\documentclass[11pt]{article}

\usepackage[a4paper,margin=1.55cm]{geometry}
\usepackage{amsmath}
\usepackage{array}
\usepackage{booktabs}
\usepackage{longtable}
\usepackage{pdflscape}
\usepackage{siunitx}
\usepackage[table]{xcolor}
\usepackage[hidelinks]{hyperref}
\usepackage{microtype}

\sisetup{detect-all}
\newcolumntype{L}[1]{>{\raggedright\arraybackslash}p{#1}}
\newcolumntype{C}[1]{>{\centering\arraybackslash}p{#1}}
\newcommand{\run}[1]{\texttt{TS#1}}
\newcommand{\file}[1]{\path{#1}}
\newcommand{\NA}{--}

\title{PMMA Rate-and-State Simulation Campaign\\
\large Parameter and strategy record from TS0117 to TS0126}
\author{Internal simulation note}
\date{1 September 2026}

\begin{document}
\maketitle

\begin{abstract}
This note records the controlled changes made in production runs \run{0117}
through \run{0126}.  It separates parameters read directly from the committed
TOML inputs from qualitative diagnoses made from the rupture maps and energy
histories.  The campaign first corrected time integration, loading amplitude,
and the rupture-stop criterion; restored the fully explicit normal-loading
history; then tested progressively slower shear loading and several
leading-edge rate-and-state profiles; and finally moved to a geometric
chamfer.  The principal conclusion is that slower displacement loading is
preferable for a physically interpretable experiment, but it is not by itself
a remedy for reverse rupture.  A rapid loading pulse may create a visually
clean front, but it changes the problem into a dynamically triggered rupture.
The recommended next design is slow prestressing followed by an immediate
fixed-displacement hold once a contiguous loading-end nucleus is detected.
\end{abstract}

\section{Scope and evidence}

The numerical values in this document were transcribed from the following
committed inputs:
\begin{quote}\small
\file{cases/rsf_0117_q4_explicit_10h.toml} through
\file{cases/rsf_0126_q4_chamfer20x5_12h.toml}.
\end{quote}
The workflow notes in \file{docs/HPC_CASES.md} and the Git history were used to
recover the purpose of each change.  Statements about forward or reverse
rupture are campaign observations from the corresponding
\file{mu_eff_map_phase_split}, animation, and energy figures; they are not
presented as new automated measurements.  At the time of writing, \run{0126}
had completed but its post-processing was still running, so its rupture
outcome is deliberately marked pending.

The standard run label is \run{0126}; ``TS00126'' is treated as a typographical
variant of the same run.

\section{Parameters held fixed}

Unless a row below states otherwise, all ten production runs used the common
values in Table~\ref{tab:fixed}.

\begin{table}[htbp]
  \centering
  \footnotesize
  \caption{Parameters held fixed across the production campaign.}
  \label{tab:fixed}
  \begin{tabular}{@{}L{5.0cm}L{5.2cm}L{4.2cm}@{}}
    \toprule
    Quantity & Value & Interpretation \\
    \midrule
    Moving block & $200\times500\ \mathrm{mm^2}$ & Loading end at $y=0$; leading end at $y=500\ \mathrm{mm}$ \\
    Stationary block & $145\times550\ \mathrm{mm^2}$ & Plane-strain two-block geometry \\
    Material & $E=7662\ \mathrm{MPa}$, $\nu=0.2$, $\rho=1.148\times10^{-9}\ \mathrm{tonne/mm^3}$ & PMMA model \\
    Normal load target & $u_n=0.6916209867\ \mathrm{mm}$ & Calibrated to approximately $16\ \mathrm{MPa}$ normal stress \\
    Element and mesh & Q4, $h=0.50\ \mathrm{mm}$ & Approximately 1.44 million displacement DOFs before the chamfer \\
    Explicit solver & \texttt{float32}, CFL $=0.35$, contact safety factor $=0.50$ & One global Tatva/JAX operator on one H200 \\
    Shear loading shape & Half cosine & Prescribed at the moving-block loading face \\
    Quasi-static shear fraction & 0 & All tangential loading is explicit dynamics \\
    Initial/reference friction & $f_0=0.8$ unless overridden locally & Low-speed reference value \\
    Reference velocities & $V_0=10^{-4}\ \mathrm{mm/s}$, $V_{\mathrm{cal}}=2000\ \mathrm{mm/s}$ & Middle zone is calibrated to $\mu=0.45$ at $V_{\mathrm{cal}}$ \\
    Shared state distance & $D_c=3.7650493\times10^{-4}\ \mathrm{mm}$ & Derived from the $X_c=5\ \mathrm{mm}$ cohesive-zone anchor \\
    Loading zone & $a=b=0.004$, length $30\ \mathrm{mm}$ & Velocity neutral, $a-b=0$ \\
    Middle zone & $a=0.005$, $b=0.0258194007$ & Velocity weakening, $a-b=-0.0208194007$ \\
    Boundary conditions & Loading face locked in shear during normal loading & Prevents tangential loading during the normal phase \\
    \bottomrule
  \end{tabular}
\end{table}

The elastic shear modulus and shear-wave transit time implied by these common
parameters are
\begin{equation}
  G=\frac{E}{2(1+\nu)}=3192.5\ \mathrm{MPa},\qquad
  c_s=\sqrt{G/\rho}=1668\ \mathrm{m/s},\qquad
  t_s=\frac{500\ \mathrm{mm}}{c_s}=0.300\ \mathrm{ms}.
  \label{eq:transit}
\end{equation}
These values are used below to distinguish slow elastic prestressing from a
rapid rupture process.

\section{Loading and numerical changes}

For a half-cosine displacement ramp of amplitude $D$ and duration $T$,
\begin{align}
  u(t) &= \frac{D}{2}\left[1-\cos\left(\frac{\pi t}{T}\right)\right],
  &0\leq t\leq T, \\
  v_{\max} &= \frac{\pi D}{2T},
  &|a(0)|=|a(T)|=\frac{\pi^2D}{2T^2}.
  \label{eq:halfcosine}
\end{align}
Table~\ref{tab:loading} includes the resulting peak boundary speed.  A
``relaxed'' normal phase means the short, damped normal-loading construction
used before \run{0120}; ``explicit'' means the restored 40-ms undamped history
with friction active and velocity carried continuously into shear.

\begin{landscape}
\begin{longtable}{@{}C{1.25cm}L{2.4cm}C{1.35cm}C{1.35cm}C{1.55cm}C{1.35cm}L{3.3cm}C{1.55cm}C{1.75cm}L{2.8cm}@{}}
  \caption{Loading, stopping, integration, and output changes. Times are in ms,
  displacement in mm, speed in mm/s, and $\Delta t$ in ns.}
  \label{tab:loading}\\
  \toprule
  Run & Normal construction & $D$ & $T_{\rm shear}$ & $v_{\max}$ & $\Delta t$ & Loading-stop criterion & Bulk shear frames & Interface shear frames & Dump changes \\
  \midrule
  \endfirsthead
  \toprule
  Run & Normal construction & $D$ & $T_{\rm shear}$ & $v_{\max}$ & $\Delta t$ & Loading-stop criterion & Bulk shear frames & Interface shear frames & Dump changes \\
  \midrule
  \endhead
  \bottomrule
  \endfoot
  \run{0117} & Relaxed: phase 3, ramp 0.5, damping time 0.2 & 2.45 & 10 ramp + 19 hold & 384.85 & 15 & One station at $y=440$: slip $\geq D_c$, $|V|\geq10$ & 36,000 & 300,000 & 1.20-TB cap; all standard bulk fields \\
  \run{0118} & Same relaxed normal phase & 2.22 & 29, no planned hold & 120.25 & 10 & Same single $y=440$ station & 36,000 & 300,000 & Same as \run{0117} \\
  \run{0119} & Same relaxed normal phase & 2.45 & 29 & 132.71 & 10 & 100\% coverage over $y=0.5$--499 plus $|V|\geq500$ & 36,000 & 300,000 & Same as \run{0118} \\
  \run{0120} & Undamped explicit: phase 40, ramp 20 & 2.45 & 29 & 132.71 & 10 & Same full-fault swept-front criterion & 36,000 & 300,000 & Sparse normal output; 1.20-TB cap \\
  \run{0121} & Same explicit normal phase & 2.45 & 43 & 89.50 & 10 & Same & 53,379 & 444,828 & Bulk strain omitted; physical shear sampling retained \\
  \run{0122} & Same explicit normal phase & 2.45 & 57 & 67.52 & 10 & Same & 70,759 & 589,655 & Bulk strain and velocity omitted \\
  \run{0123} & Same explicit normal phase & 2.45 & 75 & 51.31 & 10 & Same & 83,200 & 800,000 & Dump cap raised to 1.40 TB \\
  \run{0124} & Same explicit normal phase & 2.45 & 75 & 51.31 & 10 & Same & 83,200 & 800,000 & Same as \run{0123} \\
  \run{0125} & Same explicit normal phase & 2.45 & 75 & 51.31 & 10 & Same & 83,200 & 800,000 & Same as \run{0124} \\
  \run{0126} & Same explicit normal phase & 2.45 & 75 & 51.31 & 10 & Coverage ends at $y=479$ because contact ends at 480 & 83,200 & 800,000 & Estimated compressed dump 1.394 TB \\
\end{longtable}
\end{landscape}

Three changes in Table~\ref{tab:loading} must not be mistaken for friction-law
effects:
\begin{enumerate}
  \item \run{0118} reduced the step from 15 to 10 ns, reduced the target
  displacement from 2.45 to 2.22 mm, and spread the ramp over the full 29 ms.
  \item \run{0119} repaired a premature stopping rule.  A single far-edge
  station could be activated independently of the forward front, so the new
  rule required a contiguous full-fault sweep and raised the dynamic velocity
  threshold from 10 to 500 mm/s.
  \item \run{0120} changed the initial stress construction by restoring the
  40-ms fully explicit normal history.  It was therefore the first member of
  the later, directly comparable loading family.
\end{enumerate}

\section{Fault-profile and geometry strategies}

Table~\ref{tab:profile} records every production change to the terminal RSF
profile.  The middle zone and shared $D_c$ remained fixed.  Here $L_{\rm VS}$
is the length of the final plateau and $L_{\rm tr}$ is the half-cosine
transition immediately before it.  Until \run{0121}, the loading-end and
leading-end transition lengths were both 10 mm.  In \run{0122} the common
transition became 50 mm, so this run also smoothed the loading-end profile.
From \run{0124} onward the transitions were separated: the loading-end
transition stayed 50 mm while the leading-end transition became 100 mm.

\begin{table}[htbp]
  \centering
  \footnotesize
  \caption{Leading-edge RSF and geometry changes.  An omitted local $f_0$
  inherits the global value 0.8.}
  \label{tab:profile}
  \setlength{\tabcolsep}{4pt}
  \begin{tabular}{@{}C{1.5cm}C{0.85cm}C{0.85cm}C{1.5cm}C{1.8cm}C{1.15cm}C{1.15cm}L{5.6cm}@{}}
    \toprule
    Run & $f_{0,L}$ & $a_L$ & $b_L$ & $a_L-b_L$ & $L_{\rm VS}$ (mm) & $L_{\rm tr}$ (mm) & Intended strategy \\
    \midrule
    \run{0117}--\run{0120} & 0.80 & 0.008 & 0.005 & $+0.003$ & 30 & 10 & Baseline moderate velocity-strengthening (VS) barrier \\
    \run{0121} & 0.80 & 0.010 & 0.004 & $+0.006$ & 40 & 10 & Stronger and longer VS patch placed over the reverse nucleus near $y\simeq458$ mm \\
    \run{0122} & 0.70 & 0.012 & 0.004 & $+0.008$ & 50 & 50 & Low-$f_0$, strongly VS buffer intended to release prestress stably without unstable slip \\
    \run{0123} & 0.80 & 0.005 & 0.0258194 & $-0.0208194$ & 0 & \NA & Controlled ablation: make the terminal fault identical to the middle VW zone \\
    \run{0124} & 0.80 & 0.008 & 0.005 & $+0.003$ & 30 & 100 & Restore moderate VS but remove the sharp material-gradient scale with a long transition \\
    \run{0125} & 0.80 & 0.006 & 0.005 & $+0.001$ & 30 & 100 & Near-velocity-neutral (near-VN) terminal patch to reduce impedance while retaining weak stability \\
    \run{0126} & 0.80 & 0.008 & 0.005 & $+0.003$ & 30 & 100 & Restore \run{0124} friction; remove a $20\times5$ mm triangular wedge from the moving leading corner \\
    \bottomrule
  \end{tabular}
\end{table}

The \run{0126} chamfer is a continuous coordinate deformation of the Q4 mesh,
not a staircase deletion.  The moving block narrows linearly over the last
20 mm along the fault, reaching a 5-mm setback at $y=500$ mm.  The physical
contact ends at $y=480$ mm, whereas the RSF profile remains anchored to the
original 500-mm coordinate.  Consequently, every surviving contact station
has exactly the \run{0124} friction parameters; geometry is the only physical
intervention.

\section{Run-by-run purpose and observed response}

\begin{longtable}{@{}C{1.25cm}L{5.3cm}L{8.2cm}@{}}
  \caption{Purpose and current interpretation of each production run.}
  \label{tab:chronology}\\
  \toprule
  Run & Controlled question & Observation and interpretation \\
  \midrule
  \endfirsthead
  \toprule
  Run & Controlled question & Observation and interpretation \\
  \midrule
  \endhead
  \bottomrule
  \endfoot
  \run{0117} & Does the $X_c=5$ mm RSF calibration rupture on the target 0.5-mm Q4 mesh with a 2.45-mm, 10-ms ramp? & Rupture was reported to finish near 11 ms, leaving most of the 29-ms phase unused.  Propagation was not cleanly self-sustaining and appeared to stall in the middle.  The 15-ns step was judged too close to the stability limit.  This motivated a 10-ns step and a slower terminal ramp. \\
  \run{0118} & Can a lower 2.22-mm target and a 29-ms ramp reach full rupture with less excess loading? & A supporting 1-mm pilot nucleated dynamically at 21.208 ms, stopped the loading face at 25.051 ms and 2.120 mm, and reached the far end at 25.188 ms.  It stalled for 1.997 ms near $y=220$ mm.  Early-stop and 7-ms-hold pilots arrested near $y=137$--147 mm.  Thus the front was assisted by continuing edge loading and was not fully spontaneous.  A separate far-edge event also made the single-$y=440$ stop criterion unsafe. \\
  \run{0119} & Does a full-coverage stop rule remove false early stopping while restoring the 2.45-mm margin? & This run isolated stopping logic from the friction law.  The campaign record does not contain a separate quantitative rupture verdict; its main contribution was the robust swept-front criterion used thereafter. \\
  \run{0120} & Does restoring the original 40-ms fully explicit normal history recover the cleaner behavior of the earlier 0116 model? & Forward rupture improved, but a reverse rupture remained at the leading edge.  This was the first clear indication that terminal stress concentration, rather than only normal-state preparation or premature loading stop, controlled the unwanted second nucleus. \\
  \run{0121} & Will a 43-ms ramp plus a longer, twice-stronger VS barrier suppress the reverse nucleus? & Reverse rupture remained visible, and the energy release was less abrupt than in the desired LSW reference.  Stronger VS alone did not remove the terminal stress concentration. \\
  \run{0122} & Will a low-strength, strongly VS, broad RSF buffer creep stably and unload the terminal stress concentration? & Reverse rupture and the prolonged energy-release tail remained.  The combined $f_0$, VS amplitude, plateau length, transition length, and loading-rate changes make this result mechanistically informative but not a one-parameter test. \\
  \run{0123} & Is the VS heterogeneity itself creating the reverse rupture? & Removing the terminal VS region made reverse rupture worse.  Therefore the VS region had some suppressive value, but it was not sufficient in the earlier forms. \\
  \run{0124} & Can a moderate 30-mm VS plateau with a 100-mm transition suppress the interface-gradient stress riser? & The smoother transition still did not eliminate the leading-edge stress concentration or reverse rupture.  Attention shifted from friction-profile tuning to the physical corner geometry. \\
  \run{0125} & Does reducing the terminal patch from $a-b=+0.003$ to $+0.001$ permit enough stable release without creating a weak patch? & The result was reported as almost unchanged.  Near-VN tuning did not materially improve the reverse-rupture problem. \\
  \run{0126} & Does removing the moving-block corner stress singularity work better than further friction-law tuning? & Simulation complete; post-processing pending at the time of this note.  No rupture conclusion should be assigned until the phase-split map, energy history, and front chronology are available. \\
\end{longtable}

\section{Supporting pilot strategies}

The production sequence was supported by lower-cost pilot scans.  They should
be remembered because they constrain the interpretation of the production
runs even though they do not carry TS run numbers.

\begin{itemize}
  \item \textbf{Displacement and natural arrest.}  Target displacements around
  1.835, 1.850, 1.880, 1.900, and 2.220 mm were tested with 29-ms ramps.  The
  1.835--1.900-mm cases asked whether a front could propagate after the loading
  ramp ended naturally; the 2.220-mm case retained a far-end loading stop.
  \item \textbf{Ramp and hold.}  Ramps of 22, 25.5, and 29 ms were tested.  The
  22-ms cases included a 7-ms fixed-displacement hold.  Arrest during the hold
  showed that ending external work early can expose insufficient stored energy
  rather than merely ``cleaning'' the same rupture.
  \item \textbf{Loading-end nucleation zone.}  Velocity-neutral loading zones
  of 5 and 30 mm were examined.  Neutral direct/evolution coefficients included
  $a=b=0.002$, 0.003, 0.004, 0.005, and 0.0258194.  Additional pilots changed
  the middle direct effect while preserving the calibrated steady friction and
  fracture energy.  These tests motivated the production choice
  $a=b=0.004$ over 30 mm.
  \item \textbf{Cohesive-zone calibration.}  $X_c=3.5$, 4, and 5 mm anchors
  were bracketed, corresponding to shared RSF $D_c$ values from approximately
  $2.64\times10^{-4}$ to $3.77\times10^{-4}$ mm.  Production retained the
  5-mm anchor.
  \item \textbf{Stop position and event definition.}  Stops near $y=30$, 40,
  and 440 mm, natural no-stop cases, and eventually a full-fault contiguous
  coverage test were compared.  This progression separated true forward-front
  arrival from isolated corner or far-edge activation.
  \item \textbf{GPU scaling.}  One and two visible H200 GPUs took 268.995 and
  270.888 s in the scaling pilot.  The second GPU was idle because the current
  operator has no spatial domain decomposition; production therefore remained
  on one GPU.
\end{itemize}

\section{Is slower shear-displacement loading a good strategy?}

\subsection{What the present runs establish}

All of the imposed ramps are slow compared with bulk elastic communication.
The 10-ms ramp of \run{0117} already spans approximately 33 shear-wave transit
times, and the 75-ms ramp spans approximately 250.  The peak loading speed is
only $2.3\times10^{-4}c_s$ in \run{0117} and
$3.1\times10^{-5}c_s$ in \run{0123}--\run{0126}.  Therefore the distinction
between ``fast'' and ``slow'' in this campaign is not whether the blocks can
elastically equilibrate.  It is whether the boundary is still moving during
the rupture, how much dynamic acceleration is imposed, and how much external
work is added after nucleation.

Lengthening the 2.45-mm ramp from 29 to 75 ms reduced the peak boundary speed
from 132.7 to 51.3 mm/s (a 61.3\% reduction) and reduced endpoint acceleration
from 14.38 to 2.15 m/s$^2$ (an 85.1\% reduction).  Nevertheless, reverse
rupture persisted from \run{0120} through \run{0125}.  This is strong evidence
that loading rate is not the root cause of the reverse front.  It is not a
strict one-factor experiment, however, because the leading-edge RSF profile
also changed in \run{0121}--\run{0125}.

The \run{0118} pilot illustrates the external-work issue quantitatively.  With
the 2.22-mm, 29-ms half cosine, loading-end rupture began at 21.208 ms, when
the imposed displacement was approximately 1.847 mm.  The boundary was not
stopped until 25.051 ms and 2.120 mm.  Thus approximately 0.272 mm of additional
boundary displacement was imposed after dynamic nucleation.  Even though that
ramp was slower than \run{0117}, the rupture was still externally assisted.

\subsection{Why a deliberately fast trigger is risky}

A faster displacement ramp can make the loading-end disturbance dominate all
other nuclei and may therefore produce a visually cleaner one-way front.  That
does not necessarily demonstrate a cleaner spontaneous rupture.  It also:
\begin{itemize}
  \item increases boundary acceleration and elastic-wave content;
  \item raises instantaneous external power $P=F_{\rm reaction}\dot{u}$;
  \item can carry the front through a marginal or otherwise arresting region;
  \item makes rupture speed and energy release depend on the artificial drive;
  \item weakens comparison with a slowly driven laboratory experiment.
\end{itemize}
Such a run is useful as a diagnostic ``forced-rupture'' control, but it should
not be the primary physical case unless the experiment itself contains a
comparably rapid actuator pulse.

\subsection{Recommended loading protocol}

The recommended strategy is neither indefinite slow loading nor a rapid kick.
It is a separated preparation-and-event protocol:
\begin{enumerate}
  \item Build the normal stress with the established 40-ms fully explicit
  history.
  \item Apply the shear displacement slowly enough to keep the prestress field
  quasi-static.  The existing 75-ms half cosine is conservative; a 29--50-ms
  ramp is also quasi-static by the transit-time criterion and may reduce cost.
  \item Detect a \emph{contiguous loading-end nucleus}, not one corner node and
  not a far-edge station.  A practical criterion is a 10--20-mm-long band
  outside the known 0--6-mm corner-slip zone with slip $\geq D_c$ and dynamic
  $|V|\geq500\ \mathrm{mm/s}$.
  \item Freeze the loading-face displacement immediately when that nucleus is
  established.  Propagation after this point is powered by stored elastic
  energy, so arrest or full rupture becomes a physically meaningful outcome.
  \item Report the post-nucleation external work, the imposed displacement at
  nucleation, kinetic/stored-energy ratio, nucleation coordinate, and any
  reverse-front onset time for every run.
\end{enumerate}

This protocol directly satisfies the original requirement that the actuator
must not continue adding appreciable energy during rupture.  It also avoids
using a high-speed trigger to conceal the leading-edge stress concentration.
If the front arrests after the boundary is frozen, the correct conclusion is
that the current prestress and RSF profile do not support self-sustaining full
rupture; the remedy should then target prestress, geometry, or a physically
justified loading-end nucleation feature rather than loading faster.

\subsection{Minimal controlled loading-rate test}

The campaign does not yet contain a clean one-parameter loading-rate sweep at
fixed friction and geometry.  After \run{0126} post-processing, the least
ambiguous follow-up is a coarse-mesh pilot using the \run{0126} geometry and
\run{0124} RSF profile with $D=2.45$ mm and ramp durations 29, 50, and 75 ms.
All three should use the same contiguous loading-end stop described above.
Compare:
\begin{equation}
  y_{\rm nuc},\quad u_{\rm nuc},\quad
  \frac{E_k}{E_{\rm stored}}\bigg|_{\rm nuc},\quad
  W_{\rm ext}(t>t_{\rm nuc}),\quad
  c_R(y),\quad t_{\rm reverse}-t_{\rm nuc}.
\end{equation}
If these quantities converge as ramp duration increases, the shortest
converged ramp is the efficient choice.  If only the 29-ms case gives a clean
front and it has substantially larger post-nucleation work or kinetic energy,
that front should be classified as dynamically forced rather than selected as
the physical production solution.

\section{Campaign conclusions}

\begin{enumerate}
  \item The 10-ns step, fully explicit 40-ms normal phase, and full-fault event
  definition are numerical/workflow corrections that should remain fixed.
  \item Slowing the ramp reduced dynamic boundary forcing but did not eliminate
  reverse rupture.  The leading-edge stress concentration is the more likely
  controlling mechanism.
  \item Terminal VS has some suppressive effect: complete removal in
  \run{0123} made reverse rupture worse.  Increasing VS strength, lowering
  local $f_0$, broadening the transition, and approaching VN did not eliminate
  the problem.
  \item \run{0126} is the first production test that changes the physical
  corner geometry while restoring a previously tested RSF profile.  Its
  pending post-processing is therefore more diagnostic than another compound
  friction-profile change.
  \item Slow prestress followed by fixed-displacement rupture is preferred to
  a rapid trigger.  A fast-trigger case may still be retained as a clearly
  labelled control.
\end{enumerate}

\end{document}
